\documentclass[11pt]{article}

\usepackage[margin=1in]{geometry}
\usepackage{hyperref}

\title{Evaluating Behavioral Risk in Legal Advice Assistants}
\author{ }
\date{\today}

\begin{document}
\maketitle

\begin{abstract}
This white paper summarizes an evaluation framework for legal advice assistants focused on
behavioral risk under proposed Tennessee AI chatbot restrictions. It outlines the dataset
structure, heuristic scoring for emotional support and companionship signals, and reporting
that highlights statute-specific violations alongside legal correctness.
\end{abstract}

\section{Overview}
We assess responses to landlord-tenant prompts with three scoring axes: legal correctness,
unauthorized practice of law (UPL) risk, and Tennessee behavioral risk. The behavioral axis
maps to four clause categories: emotional support, companion framing, friendship simulation,
and human simulation. The framework is designed to surface behavioral risks even when legal
content is accurate.

\section{Method}
Prompts are grouped by scenario and include neutral, anxious, and support-request variants.
Each run logs the full conversation, a final answer, and per-axis scores. Clause-specific
hits are reported separately to maintain interpretability.

\section{Implications}
Results typically show weak correlation between legal correctness and behavioral compliance,
indicating the need for separate guardrails. The approach can be extended to other domains
where the legal risk hinges on relational framing rather than factual accuracy.

\section{Limitations}
The scoring is heuristic and intentionally conservative; human review is recommended for
edge cases. Dataset coverage is limited to landlord-tenant scenarios and should be expanded
for broader claims.

\section{Next Steps}
Future work includes multi-model comparisons, calibrated thresholds for conversational
markers, and a larger dataset with jurisdictional diversity.

\end{document}
